\documentclass[12pt]{article}
%%%%%%%%%%%%%%%%%%%%%%%%%%%%%%%%%%%%%%%%%%%%%%%%%%%%%%%%%%%%%%%%%%%%%%%%%%
\usepackage{amssymb}
\usepackage{amsfonts} 
\usepackage{amsmath}
\usepackage{natbib}
\usepackage{hyperref}
\usepackage{mathabx} %to get widebar
\hypersetup{
  colorlinks,
  citecolor=blue,
  linkcolor=red,
  urlcolor=blue}
\usepackage{rotating} %to rotate text 90 degrees
\usepackage{geometry}
\usepackage{graphicx}
\usepackage{setspace}
\usepackage{comment}
\usepackage{float}
\usepackage{dcolumn}
\usepackage{setspace}
\usepackage[flushleft]{threeparttable}
\usepackage[symbol]{footmisc}
\usepackage{url}
\usepackage{subcaption}
\usepackage{graphicx}
\usepackage{multirow}
\usepackage{color, colortbl}
\usepackage{bbm}
\usepackage{lscape}
\definecolor{Gray}{gray}{0.9}
\renewcommand{\thefootnote}{\fnsymbol{footnote}}
\setcounter{MaxMatrixCols}{10}
\newtheorem{theorem}{Theorem}
\newtheorem{acknowledgement}[theorem]{Acknowledgement}
\newtheorem{algorithm}[theorem]{Algorithm}
\newtheorem{axiom}[theorem]{Axiom}
\newtheorem{case}[theorem]{Case}
\newtheorem{claim}[theorem]{Claim}
\newtheorem{conclusion}[theorem]{Conclusion}
\newtheorem{condition}[theorem]{Condition}
\newtheorem{conjecture}[theorem]{Conjecture}
\newtheorem{corollary}[theorem]{Corollary}
\newtheorem{criterion}[theorem]{Criterion}
\newtheorem{definition}[theorem]{Definition}
\newtheorem{example}[theorem]{Example}
\newtheorem{exercise}[theorem]{Exercise}
\newtheorem{lemma}[theorem]{Lemma}
\newtheorem{notation}[theorem]{Notation}
\newtheorem{problem}[theorem]{Problem}
\newtheorem{proposition}[theorem]{Proposition}
\newtheorem{remark}[theorem]{Remark}
\newtheorem{solution}[theorem]{Solution}
\newtheorem{summary}[theorem]{Summary}
\newenvironment{proof}[1][Proof]{\noindent\textbf{#1.} }{\ \rule{0.5em}{0.5em}}
\geometry{top=1.2in,bottom=1.2in,left=1.2in,right=1.2in}
%\geometry{top=1.425in,bottom=1.425in,left=1.5in,right=1.5in}
%\geometry{top=1.5in,bottom=1.5in,left=1.35in,right=1.35in}
\renewcommand{\rmdefault}{ptm}
\setcounter{topnumber}{2}
\setcounter{bottomnumber}{2}
\setcounter{totalnumber}{4}
\renewcommand{\topfraction}{0.85}
\renewcommand{\bottomfraction}{0.85}
\renewcommand{\textfraction}{0.15}
\renewcommand{\floatpagefraction}{0.7}
\usepackage{setspace} 
\renewcommand{\thefootnote}{\arabic{footnote}}
\setlength{\bibsep}{3.9pt}
\linespread{1.5}
\usepackage{chngcntr}
\usepackage{etoolbox}
\makeatletter
\AtBeginDocument{\preto\@title{\singlespacing}}% in case it isn't single spaced
\makeatother
\usepackage{tikz} %to draw figures in latex
\usepackage{ragged2e}

\begin{document}


\title{{\bf \Large Deep in the Tails: Multi-Task Learning for Macroeconomic Tail Risk}\thanks{The views expressed in this paper are those of the authors and do not necessarily reflect those of the Bank of Canada. We thank ...}}

%Alternative titles:
%Deep in the Tails: Multi-Task Learning for Macroeconomic Tail Risk
%Don't be mean, get into shape! Multi-task neural network for macroeconomic risks

\author{Robert \textsc{Proner}\thanks{University of Toronto. Email: \texttt{robert.proner@mail.utoronto.ca}.}} 
% \qquad Thibaut \textsc{Duprey}\thanks{Bank of Canada, 234 Wellington Street West, Ottawa, ON K1A0G9, Canada. Email: \texttt{tduprey@bank-banque-canada.ca}. Corresponding author.} \qquad Fred \textsc{Liu}\thanks{University of Guelph. Email: \texttt{fred.liu@uoguelph.ca}.}}

\date{\today}
%\date{}
\maketitle
\begin{abstract}
\noindent 
We develop a deep multi-task neural network (DMQ) to forecast the conditional distribution of inflation, industrial production, and unemployment from 1998-2024. As benchmarks, we use quantile regressions from the latest macroeconomic risk literature and a large set of popular machine-learning methods. We find DMQ provides substantial gains in forecast accuracy for the conditional distribution as well as the conditional mean, especially during the GFC and Covid.   
\end{abstract}
\textbf{Keywords}: \\
\textbf{JEL Codes}: 


\section{Introduction}

In recent years policy makers have become increasinly concerned with uncertainty surrounding macroeconomic forecasts. 
For example, the Federal Open Market Committee (FOMC) regularly includes distributions around its predictions for 
key macroeconomic indicators within its outlook publications. Distributions around forecasts are important because
point forecasts fail to paint an accurate picture of the level of risk in the economy. 
Policy makers are particularly concerned with tail risks and economic uncertainty clearly is factored in their decisions
\citep{andradeTAILSINFLATIONFORECASTS2012}. It is therefore important to have accurate forecasts of the 
conditional distribution of important macroeconomic indicators. 

Forecasting the conditional distribution of macroeconoic indicators is challenging for several reasons. First, 
the conditional distribution is highly time-varying, both in location, scale, and shape \citep{adrianVulnerableGrowth2019,lopez-salidoInflationRisk2024,kileyUnemploymentRisk2022,clarkRealTimeDensityForecasts2011,pruserNonlinearitiesMacroeconomicTail2024,andradeTAILSINFLATIONFORECASTS2012}.
Second, the functional form linking the quantiles of the conditional distribution to the predictors is likely 
non-linear \citep{adrianVulnerableGrowth2019,lopez-salidoInflationRisk2024,pruserNonlinearitiesMacroeconomicTail2024},
and depends on a large number of macroeconomic and financial predictors \citep{pruserNonlinearitiesMacroeconomicTail2024,carriero_specification_2025,adrianMachinelearningGrowthRisk2025}.
Third, linear models, specifically quantile regressions (QR) are susceptible to overfitting and quantile crossing \citep{schmidt_quantile_nodate}.
In this paper we remedy these issues by developing a deep multi-quantile neural network (DMQ) to forecast the conditional distribution 
of inflation, industrial productions (IP) growth, and percent change in unemployment. We demonstrate 
that DMQ provides significant gains in out-of-sample forecast accuracy in both the conditional quantiles 
as well as the conditional mean, especially during the GFC and Covid. We find [INSERT RESULTS ABOUT FEATURE IMPORTANCE].
We also show that forecasting the entire conditional distribution ointly eliminates the need for quantile-crossing
corrections and provides more accurate forecasts than forecasting each quantile separately.



\section{Method}\label{section:data}
\subsection{...}





%\clearpage

\section{Conclusion} \label{section:conclusion}



\newpage
\singlespacing

%%% BibTeX style.
\bibliographystyle{ecta}%{aer}
%% BibTeX database file.
\bibliography{references.bib}

\clearpage
\appendix
\counterwithin{figure}{section}
\counterwithin{table}{section}


\appendix

\section{Appendix}



\end{document}